\documentclass[11pt,a4paper]{article}
\usepackage[margin=27mm]{geometry}
\usepackage{amsmath,amssymb}
\usepackage[T1]{fontenc}
\setlength{\emergencystretch}{2em}
\title{An analytic spurious stationary point with a positive-measure basin\\
\large Disproof of conjecture 00000006073}
\author{ziangni-sys}
\date{4 October 2026}
\begin{document}
\maketitle
The source explicitly defines spurious minima as stationary
points with zero gradient that are not minima. Its final
conjunct states that their attraction basins have measure zero,
including in the analytic case.
The analytic function $f(x)=x^3$ disproves that assertion for
the standard continuous gradient flow.
The separate smoothing and Morse-index assertions are unused.

\textbf{Actual analytic objective and stationary point.}
Work on the real Euclidean line and set
\[
 f(x)=x^3.
\]
This polynomial is real analytic at every point. Its actual
derivative and Hilbert-space gradient are
\[
 f'(x)=\nabla f(x)=3x^2.
\]
In particular $\nabla f(0)=0$.

The origin is not a local minimum. Every neighborhood contains
a negative point, where $f(x)<0=f(0)$.
More explicitly, if a neighborhood ball has radius $\varepsilon>0$,
then $-\varepsilon/2$ lies in that ball and satisfies
\[
 f(-\varepsilon/2)=-(\varepsilon/2)^3<0.
\]
Thus the origin is a spurious stationary point under the
source's explicit definition.
It is a degenerate stationary inflection point, rather than a
nondegenerate strict saddle. No such nondegeneracy hypothesis
appears in the source.

\textbf{Exact gradient flow for all positive starts.}
For each $s>0$, define
\[
 u_s(t)=\frac{s}{1+3st},\qquad t\geq0 .
\]
The denominator is strictly positive for every forward time,
and $u_s(0)=s$. Direct differentiation gives
\[
 u_s'(t)=-\frac{3s^2}{(1+3st)^2}
        =-3u_s(t)^2=-\nabla f(u_s(t)).
\]
This is an actual global forward gradient-flow solution.

\newpage
\textbf{Convergence of every such trajectory.}
Since $s>0$, the denominator tends to infinity as $t\to\infty$.
Consequently
\[
 \lim_{t\to\infty}u_s(t)=0.
\]
All positive initial conditions therefore lie in the attraction
basin of the spurious point.

For clarity, define the forward-flow basin directly by
\[
 B=\left\{s\in\mathbb R:
 \begin{array}{l}
 \text{there exists a forward trajectory }u,\ u(0)=s,\\
 u'(t)=-\nabla f(u(t))\text{ for all }t\geq0,\\
 u(t)\longrightarrow0\text{ as }t\to\infty
 \end{array}\right\}.
\]
The derivative in this definition is the genuine real
derivative, not a prescribed scalar value. The vector field
$-3x^2$ is smooth and locally Lipschitz, so its usual local
uniqueness also agrees with the standard gradient-flow basin.
For the counterexample only the explicit global trajectories
above are needed.

\textbf{Actual positive Lebesgue measure.}
The preceding formulas prove
\[
 (0,1)\subseteq B.
\]
For ordinary Lebesgue measure $\lambda$ on $\mathbb R$,
monotonicity and the exact interval-measure formula give
\[
 \lambda(B)\geq\lambda((0,1))=1>0.
\]
The formal proof uses the actual Lebesgue outer measure on sets,
so no measurability assumption about the basin is needed to
exclude nullity. In particular the basin cannot have measure zero.

\textbf{Scope.}
The disproof targets the attraction-basin nullity assertion,
under continuous gradient descent on the actual analytic
objective. It does not assert a failure of a theorem restricted
to nondegenerate strict saddles. The source's wider definition
expressly includes stationary nonminima such as the analytic
inflection point used here.
No claim about the separate smoothing, index or critical-point
count formulas is required.

\textbf{Formal correspondence and reproduction.}
Lean proves real analyticity, the actual derivative and gradient,
stationarity, failure of local minimality, the explicit initial
condition, the genuine differential equation at every
nonnegative time, the actual limit, interval inclusion in the
defined gradient-flow basin, and its actual Lebesgue measure
lower bound. The final theorem combines those properties.

Run \texttt{lake build} in \texttt{lean/} with Lean 4.19.0 and
public Mathlib revision:
\begin{center}\small
\texttt{c44e0c8ee63ca166450922a373c7409c5d26b00b}
\end{center}
The project prints the relevant axiom audits.
\end{document}
